\subsection{有序元素序列的复合}

回忆一下，集合 E 的元素族是从一个集合 I（称为指标集）到 E 的映射 \(t \mapsto x_{t}\) ；族 \((x_{t})_{t \in I}\) 若其指标集是有限的，则称为有限族。

一个有限族 \((x_{t})_{t\in I}\) 若其指标集 I 是全序的，则称为 E 的有序元素序列。

特别地，每个有限序列 \((x_{i})_{i\in\mathbb{H}}\) ，其中 H 是自然数集 N 的一个有限子集，若 H 被赋予由自然数之间的关系 \(m \leqslant n\) 所诱导的序关系，则可将其视为一个有序序列。

两个有序序列 \((x_{i})_{i\in\mathbf{I}}\) 和 \((y_{k})_{k\in\mathbf{K}}\) 称为相似的，如果存在一个从 I 到 K 的有序集合同构 \(\phi\) ，使得 \(y_{\phi(i)} = x_{i}\) 对所有 \(i \in I\) 成立.

每个有序序列 \((x_{\alpha})_{\alpha \in \mathbf{A}}\) 都相似于某个适当的有限序列。因为存在一个从 \(\mathbf{A}\) 到区间 \([0, n]\) （在 \(\mathbf{N}\) 中）上的递增双射.

定义4. 设 \((x_{\alpha})_{\alpha \in \mathbf{A}}\) 为广群 \(\mathbf{E}\) 中元素的一个有序序列，其指标集 \(\mathbf{A}\) 非空。在律 \(\top\) 下的有序序列 \((x_{\alpha})_{\alpha \in \mathbf{A}}\) 的合成，记作 \(\bigcap_{\alpha \in \mathbf{A}} x_{\alpha}\) ，是 \(\mathbf{E}\) 的元素，通过对 \(\mathbf{A}\) 中元素个数进行归纳定义如下：

\begin{enumerate}
\def\labelenumi{(\arabic{enumi})}
\item
  若 \(\mathbf{A} = \{\beta\}\) ，则 \(\bigcap_{\alpha \in \mathbf{A}} x_{\alpha} = x_{\beta}\) ;
\item
  若 A 有 \(p > 1\) 个元素， \(\beta\) 是 A 的最小元素，且 \(A' = A - \{\beta\}\) ，则 \(\bigcap_{\alpha \in A} x_{\alpha} = x_{\beta} \top \left( \bigcap_{\alpha \in A} x_{\alpha} \right)\) .
\end{enumerate}

由此立即得出（通过对指标集中元素个数进行归纳）两个相似有序序列的复合相等；特别地，任何有序序列的复合等于某个有限序列的复合。若 \(\mathbf{A} = \{\lambda, \mu, \nu\}\) 有三个元素（ \(\lambda < \mu < \nu\) ），则复合 \(\bigcap_{\alpha \in \mathbf{A}} x_{\alpha}\) 是 \(x_{\lambda} \top (x_{\mu} \top x_{\nu})\) .

注：注意，有序序列合成的定义存在一定的任意性；我们引入的归纳是“从右到左”进行的。如果我们“从左到右”进行，上述有序序列 \((x_{\lambda}, x_{\mu}, x_{\nu})\) 的合成将是 \((x_{\lambda} \top x_{\mu}) \top x_{\nu}\) .

\[
\left(x _ {\alpha}\right) _ {\alpha \in A}
\]

\[
\begin{array}{c} \underline {{\mathsf {I}}} \\ \alpha \in \mathbf {A} \end{array}
\]

记作 \(\sum_{\alpha \in \mathbf{A}} x_{\alpha}\) ，称为有序序列 \((x_{\alpha})_{\alpha \in \mathbf{A}}\) 的和（ \(x_{\alpha}\) 称为和的各项）；对于以乘法记法书写的律，通常记作 \(\prod_{\alpha \in A} x_{\alpha}\) ，称为有序序列 \((x_{\alpha})\) 的积（ \(x_{\alpha}\) 称为积的因子）。†

当索引集（及其顺序）不会引起混淆时，在有序序列的合成记号中通常可以省略它，例如对于以加法记法书写的律，我们可写成 \(\sum_{\alpha} x_{\alpha}\) 而不是 \(\sum_{\alpha \in A} x_{\alpha}\) ；其他记号同理。

对于用 \(\top\) 表示的律，序列 \((x_{i})\) （其指标集为非空区间 \((p,q)\) ，在 \(\mathbf{N}\) 中）的合成记作 \(\bigcap_{p\leqslant i\leqslant q}x_{i}\) 或 \(\bigcap_{i = p}^{q}x_{i}\) ；对用其他符号表示的律亦然。

设 E 和 F 是两个广群，其合成律记为 T，f 是 E 到 F 的一个同态。对于每个有序序列 \((x_{\alpha})_{\alpha \in \mathbf{A}}\) （元素属于 E ），有

\[
f \left(\underset {\alpha \in A} {T}\right) = \underset {\alpha \in A} {T} f (x _ {\alpha}).\tag{2}
\]

